% Pillar-7 (ZVF theory + adaptive-G controller) standalone abstract
\begin{abstract}
Our companion Pillar-2 paper characterized the Zero-Variance Fraction
(ZVF)---the share of GRPO groups whose completions all receive the same
reward---as a deliberately \emph{descriptive} diagnostic of signal starvation.
This paper is a retrospective controller audit and a prospective intervention
plan: we ask which ZVF-derived quantities are mechanistically predictable and
whether they are ready to control rollout allocation. This does not
overturn Pillar-2's restraint---we \emph{test} ZVF's predictive value through
falsifiable predictions (T1--T3 below) rather than assume it, and we flag where
that value rests on the weakest evidence (the causal gradient link is toy-scale
and directional only). We model the
usable learning signal per prompt as $S = p(1-p)\,(1 - h_G(p))$, where $p$ is
the latent success probability and $h_G(p) = p^G + (1-p)^G$ is the probability
a size-$G$ group is degenerate, and derive three testable consequences:
(T1) the expected ZVF indicator equals $p^G + (1-p)^G$; (T2) larger group
size $G$ lowers ZVF; (T3) gradient magnitude tracks $p(1-p)$. Each prediction
is paired with an empirical check: a 368-run Weights~\&~Biases audit across seven
models (all logged runs---a superset of the benchmark's 70+ core training runs; the
core set is defined by run \emph{type}, i.e.\ completed GRPO training runs, and is
\emph{not} filtered on convergence or reward, while the remaining logged runs are
evaluations, ablation probes, and short/aborted infrastructure runs---so the audit
is not subject to outcome-based survivorship selection)
reproduces the predicted U-shape of ZVF in accuracy (ZVF $0.95$--$0.97$ at both
reward extremes, $0.25$--$0.29$ at mid-range rewards of $0.35$--$0.50$); a
model~$\times$~$G$ grid confirms the monotone $G$ effect at interior accuracy
(e.g.\ Qwen3-32B: $0.33 \to 0.18 \to
0.08$ for $G = 4/8/16$); and a toy-scale open-backward-pass experiment finds
$\mathrm{corr}(\|\nabla\|, p(1-p)) = +0.71$ (0.5B model, synthetic arithmetic;
directional only). Because ZVF aliases mastery with incapacity, we promote the
pairwise-contrast density $\mathrm{PCD} = \frac{G-1}{G}\,\E[p(1-p)]$ as the
control signal: a micro-jitter falsification collapses batch ZVF from $0.158$
to $0.000$ while leaving PCD unchanged. We then operationalize the theory as a
\texttt{zvf-triage} callback and an adaptive-$G$ controller that escalates
group size on ZVF spikes. In a four-arm open-trainer audit, adaptive-$G$
matches the best fixed-recipe held-out gain ($+0.575$) at $186$ rollouts,
while DAPO-style dynamic sampling drives ZVF to $0.00$ at $+45\%$ rollout
cost---an accuracy-versus-compute trade, not a controller win. In the larger
2,560-observation controller trace, 1,723 of 1,867 escalation events (92.3\%)
occur on all-correct groups and only 144 on all-wrong groups. This asymmetry
shows that a symmetric ZVF trigger mostly spends on solved prompts; deleting
those fires is only a counterfactual proposal because performance preservation
has not been measured. No result here establishes prospective superiority over
static $G{=}16$, Dr.~GRPO, or dynamic sampling. All intervention evidence is
single-task, small-$n$, or retrospective, and the paper ends with the matched-
budget, multi-seed test required to earn control authority.
\end{abstract}
